\chapter{Conclusion}
\label{ch:conclusion}

This dissertation set out to test a simple but consequential question on a
controlled traffic benchmark: does an explicitly cooperative multi-agent algorithm
actually outperform a set of independent learners when the two are matched in every
other respect? The investigation produced a clear, and initially surprising, answer.

A MAPPO traffic-signal controller was implemented on a \grid{2} grid, trained to
stable convergence, and shown --- under a \emph{fair}, matched-clearance evaluation ---
to outperform both the fixed-time and max-pressure heuristics decisively. A
methodological contribution emerged along the way: the heuristic baselines had been
flattered by an unfair zero-clearance evaluation, and correcting it (charging every
controller the same safety clearance) both restored a fair comparison and
strengthened the case for the learned controllers. The project's MAPPO was further
confirmed to be a faithful implementation by matching the canonical published
recipe.

The central result is that, against this strong and faithful MAPPO, an Independent
PPO comparator --- differing only in its decentralized critic and purely local
reward --- proved \emph{statistically indistinguishable} on every traffic-performance
metric. The sophisticated coordination machinery bought no measurable advantage; its
only distinguishing behaviour was a slightly higher rate of phase switching that
produced no benefit. On a \grid{2} network, the \emph{coordination value} is
approximately zero.

The interpretation is not that coordination is worthless, but that the testbed is
too small to require it: at this scale incentives are already locally aligned, and
independent learners recover coordinated behaviour implicitly. This precisely
locates where explicit coordination should begin to matter --- at larger network
scales --- and motivates the future work of scaling to \grid{5} and of extending the
study to genuine social dilemmas, where reward structure and the temptation to
defect are expected to make coordination indispensable. The contribution of this
work is thus both a careful negative result and a sharpened hypothesis: the value of
cooperation in multi-agent reinforcement learning is not a constant but a function
of scale, and the controlled traffic setting establishes the baseline against which
that function can now be measured.
